\section{What the Composer Wrote}
\label{sec:ch07-what-the-composer-wrote}

\index{router execution config!what the command produces|(}%
Compose again and the tool says what it made:

\begin{minted}[fontsize=\footnotesize,breakanywhere]{text}
Router execution config successfully written to ".../graph/router.json".
\end{minted}

\index{router execution config!is WunderGraph's term}%
Note the noun, because it is not the one
chapter~\ref{ch:federation-model} used. That chapter called the composed
result a supergraph, which is Apollo's word for it. WunderGraph's own reference
for this command calls the output a router execution
config~\autocite{cosmocompose}, and the word supergraph does not appear on
that page at all. The two names are for different things and both are in play:
a supergraph is the arrangement, and this file is one composer's compiled form
of it, for one router.

\index{router execution config!no user documentation defines it}%
What is in the file is not documented anywhere I could find. There is no
reference page for its shape, no JSON schema published beside it. What defines
it is a Protocol Buffers message in the Cosmo source tree, and that message is
short enough to read~\autocite{cosmonodeproto}:

\begin{minted}{text}
message RouterConfig {
  EngineConfiguration engine_config = 1;
  string version = 2;
  repeated Subgraph subgraphs = 3;
  optional FeatureFlagRouterExecutionConfigs feature_flag_configs = 4;
  string compatibility_version = 5;
}
\end{minted}

\index{router execution config!five keys}%
Five fields, and the file this graph produced has exactly five keys with those
names. That is worth doing once for any artifact a tool hands you and expects
you to deploy. Seventeen kilobytes of JSON for a single-service graph looks
opaque and is not.

\subsection{The Three Small Ones}

\index{compatibilityVersion@\code{compatibilityVersion}!names the composer}%
\code{compatibilityVersion} is \code{1:0.63.3}, and the second half of that
string is a version number of something you never installed. It matches the
version of the composition package inside the tool, which I checked against
that package's own manifest. So the config appears to record which composer
produced it, and that is the only place it is recorded, which matters the day
somebody hands you a config and asks why the router will not load it.

\index{version@\code{version}!is nil without a control plane}%
\code{version} is a nil uuid, thirty-two zeroes. Composing locally is the mode
this book uses and the mode WunderGraph's own page recommends against for
production, in favor of fetching a config from their platform. What a config
fetched from that platform carries in this field, I have not seen. This one
carries nothing.

\index{subgraphs@\code{subgraphs}!carries the routing url}%
\code{subgraphs} is the list, and it holds what \code{graph.yaml} said: a
generated id, the name, and the routing url the composer never called. Here is
where that url ends up, and the router is the thing that will use it.

\subsection{Two Schemas, Not One}

\index{graphqlSchema@\code{graphqlSchema}!is what a client reads}%
\code{engineConfig} holds the rest, and the first interesting thing in it is
\code{graphqlSchema}. That is the document a client gets, and
chapter~\ref{ch:federation-model} promised what it would look like: the
directives do their work at composition and are gone from what the client
reads. They are. No \code{@key}, no \code{@shareable}, no \code{@link}, no
\code{\_entities}, no \code{\_service}, no \code{\_Any} and no \code{\_Entity}.
A client reading this schema cannot tell that federation was involved.

\index{FieldSet@\code{FieldSet}!survives into the client schema}%
Almost. One line survives that I would not have predicted:

\begin{graphqlsdl}
scalar FieldSet
\end{graphqlsdl}

\index{FieldSet@\code{FieldSet}!nothing references it}%
\code{FieldSet} is the scalar that \code{@key}'s argument is typed with, and
the directive that needed it is gone. The name appears exactly once in the
whole client schema, which is that declaration, so nothing refers to it. It is
a scalar a client can neither produce nor receive, left standing because
nothing swept up after the directive it belonged to. Harmless, and a good
reminder that the composed document is generated rather than curated.

\index{serviceSdl@\code{serviceSdl}!the subgraph document is kept whole}%
Then read one level down, into the datasource entry, and the seam is back.
Inside \code{customGraphql} there is a \code{federation} block, and inside that
a \code{serviceSdl} string holding your subgraph's own document, whole. The
\code{@key} on both types, the \code{@shareable} on both root fields,
\code{\_entities} and the \code{@link} header are all still in it. So the config holds
both views at once: the schema the client is allowed to see, and the schema the
router needs in order to serve it. Chapter~\ref{ch:blame-routing} is about what
it means that only one of those two ever reaches the outside.

\subsection{The Seam, Compiled}

\index{keys@\code{keys}!the key directives as a table}%
The last piece is the one that makes the whole file legible. Beside that
subgraph document, the composer writes out what it worked out from it:

\begin{minted}{json}
"keys": [
  {"typeName": "Session", "selectionSet": "id"},
  {"typeName": "Speaker", "selectionSet": "id"}
]
\end{minted}

\index{keys@\code{keys}!is a readout of which types are entities}%
Two \code{@key} directives, compiled into a lookup table. This is the same
readout chapter~\ref{ch:first-subgraph} found in the \code{\_Entity} union, in
the form the router actually consumes: a type name and the fields that identify
one of them. A type missing from this table is a type the router has no way
into, and later in this chapter that stops being an abstract worry.

\index{rootNodes@\code{rootNodes}!which fields this service can answer}%
Above it sit \code{rootNodes} and \code{childNodes}, which between them say
which fields this datasource can answer. Four types are root nodes, each with
its field list: \code{Query}, \code{Mutation}, \code{Session} and
\code{Speaker}. Eight are child nodes, reachable only underneath something
else, and they are every other type in the schema: \code{PageInfo}, the two
connection types, the mutation payload, the two error types, the \code{Error}
interface and \code{Node}. With one subgraph that partition is not telling you
much. With three it
is the map the router plans against, and
section~\ref{sec:ch07-the-walk} is about what the composer does with it before
it hands it over.
\index{router execution config!what the command produces|)}%
